% ============================================================
%  ABSTRACT
%  Per Prof. [RK1]: include work done AND specific results
% ============================================================
\chapter*{Abstract}
\addcontentsline{toc}{chapter}{Abstract}

Supply chain management in urban logistics networks faces a fundamental challenge:
delivery fleets must satisfy multiple competing objectives simultaneously, including
minimizing fuel costs, preserving cargo freshness, maximizing delivery speed, and
avoiding unsafe road conditions. Traditional route planning algorithms such as A*
use a single fixed cost function and cannot adapt dynamically to real-time events
like road accidents, extreme weather, or sudden demand spikes. This project addresses
this limitation by designing and implementing a Supply Chain Digital Twin for the
city of Nagpur, validated over a full 365-day simulation.

The system is built around a three-tier AI architecture. The first tier is a Central
PPO (Proximal Policy Optimization) Brain, trained for 50,000 steps to select
macro-level routing strategies from three options: Fuel-Priority, RSL-Priority, and
Speed-Priority. The second tier is a per-truck Edge Q-Learning Agent that performs
real-time emergency micro-rerouting when accidents or weather events block the
planned route. The third tier is an online Ridge Regression ETA Forecaster that
predicts segment travel times using an 8-dimensional spatial-temporal feature vector
and updates itself continuously throughout the simulation.

The system was evaluated through a 365-day comparative benchmark against a
heuristic A* baseline. The AI-enabled configuration fulfilled 323 additional orders,
delivered 474,708 more kilograms of cargo, and reduced retailer stockout events from
70,898 to 33,401, a reduction of 52.9\%. Road accidents were reduced by 135 events
as an emergent consequence of the ETA Forecaster learning to route trucks away from
high-risk zones during peak hours. Cargo freshness at delivery and fleet fuel
efficiency were identical in both configurations, at 99.7\% and 78.7\% respectively,
confirming that the throughput gains were achieved without any trade-off in quality
or operational cost. The results demonstrate Pareto optimality: six of eight
performance metrics improved simultaneously, and no metric regressed.
